\begin{titlepage}
\begin{abstract}
%
Game Software Engineering (GSE) differs from traditional enterprise software because player preferences concerning narrative, mechanics, or balance are abstract and difficult to quantify, especially prior to release. While the rise of live-service games and early-access models has shifted production toward continuous iteration, how informal player feedback can be systematically converted into formal requirements remains poorly defined. Existing analytics tools model what players do through telemetry, but not why they do it, and LLM-based feedback pipelines are documented only for mobile app stores, with no publicly documented equivalent for PC and console games.

To address this gap, this thesis proposes StARAPTOR (a Steam Automated Review Association Miner for Player-count Time-series Optimized Regression), a prototype tool that leverages Large Language Models (LLMs) to extract discrete attributes from Steam reviews. The tool expands feature dimensionality using FP-Growth association mining to identify recurring relationships, combining them into a feature matrix that feeds an XGBoost time-series regressor to forecast daily mean concurrent player counts. By mathematically linking player sentiment to retention-related metrics, StARAPTOR automates feedback analysis and helps developers quantify the impact of user requirements on player engagement.

A preliminary validation on one year of data from a single high-traffic game (Cyberpunk 2077) shows consistent extraction across models on the most consumed aspects, 73.7\% actionable mined rules, and a recursive 30-day forecast degraded by up to 51.7\% against the KPI-only baseline, though with underestimation cut by 69.6\%. Overall, the work closes the identified gap by providing a reproducible, model-agnostic pipeline that converts review feedback into quantified, traceable forecasting inputs, while showing that their forecasting value must be validated per game.
%
\\[1cm]
\end{abstract}
\end{titlepage}